% Seminar 4 - Page 2. Fragment for
% % Seminar 4 - Page 2. Fragment for
% % Seminar 4 - Page 2. Fragment for
% % Seminar 4 - Page 2. Fragment for
% \input{seminars/pages/seminar4-pagina2-en}; no preamble or forced page break.
% EDITORIAL NOTES:
% - Exercise 6 has two parts. The proof distinguishes an existing monoid
%   from a semigroup without an identity, then adjoins a new symbol, infinity.
% - The source explicitly writes x*infinity = infinity*x = x. It does not
%   separately write infinity*infinity = infinity; that case is required
%   to make the new operation a monoid operation.
% - The source first writes Z_n generically. The subsequent equality
%   3-hat*2-hat = 1-hat*2-hat, with 3-hat != 1-hat, works for n=4;
%   it is NOT true for arbitrary n. We retain Z_n and flag this issue.
% - The final matrix-monoid example is recorded as an observation, not
%   expanded into a proof: AB=CB with A != C.
% - Handwritten abbreviations and very faint show-through are not
%   reconstructed beyond what the mathematical lines support.

\subsection*{Seminar 4 - Page 2}

\noindent\underline{Exercise 6.}

\noindent\textbf{(a)} Show that every semigroup $S$ can be embedded in a
monoid: there exist a monoid $M$ and an injective semigroup
homomorphism
\[
  f:S\longrightarrow M, \qquad f(xy)=f(x)f(y).
\]

\noindent\textbf{(b)} There also exists a monoid that cannot be embedded
in a group: a monoid $M$ for which there are no group $G$ and
injective monoid homomorphism $f:M\to G$.

\medskip
\noindent\underline{Proof.}

\noindent\textbf{(a)} Let $S$ be a semigroup.

\smallskip
\noindent\textbf{I.} If $S$ is already a monoid, take $M=S$ and
\[
  i:S\longrightarrow M=S,\qquad i(x)=x.
\]

\noindent\textbf{II.} Suppose $S$ is a semigroup but not a monoid.
Take $\infty\notin S$ and put
\[
  M\overset{\mathrm{not}}{=}S\cup\{\infty\}.
\]
For $x,y\in M$ the operation is defined, as written in the notes, by
\[
\begin{gathered}
  x\circ y\ \overset{\mathrm{def}}{=}\ xy
     \quad(x,y\in S),\\
  \infty\circ x\ \overset{\mathrm{def}}{=}\
     x\circ\infty\ \overset{\mathrm{def}}{=}\ x
     \quad(x\in S).
\end{gathered}
\]
Then $(M,\circ,\infty)$ is a monoid. Consider
\[
  f:S\longrightarrow S\cup\{\infty\}=M,\qquad f(x)=x;
\]
thus $S$ is embedded in the monoid $M$.

\medskip
\noindent\textbf{(b)} The notes consider the monoid
\[
  \mathbb Z_n=\{\widehat{0},\widehat{1},\ldots,
      \widehat{n-1}\},\qquad (\mathbb Z_n,\cdot).
\]
Suppose that there exist a group $G$ and an injective monoid
homomorphism $f:\mathbb Z_n\to G$. Thus
\[
  f(xy)=f(x)f(y),\qquad f(\widehat{1})=e.
\]
The handwritten cancellation argument is
\[
\begin{aligned}
  f(\widehat{3})f(\widehat{2})
    &=f(\widehat{1})f(\widehat{2})\\
  &\Longrightarrow f(\widehat{3})=f(\widehat{1})\\
  &\overset{f\,\mathrm{injective}}{\Longrightarrow}
      \widehat{3}=\widehat{1}\qquad\text{(false).}
\end{aligned}
\]

\noindent Another example written at the bottom is the matrix monoid
$(M_n(K),\cdot)$, with matrices $A,B,C\in M_n(K)$ such that
\[
  AB=CB,\qquad A\neq C.
\]
; no preamble or forced page break.
% EDITORIAL NOTES:
% - Exercise 6 has two parts. The proof distinguishes an existing monoid
%   from a semigroup without an identity, then adjoins a new symbol, infinity.
% - The source explicitly writes x*infinity = infinity*x = x. It does not
%   separately write infinity*infinity = infinity; that case is required
%   to make the new operation a monoid operation.
% - The source first writes Z_n generically. The subsequent equality
%   3-hat*2-hat = 1-hat*2-hat, with 3-hat != 1-hat, works for n=4;
%   it is NOT true for arbitrary n. We retain Z_n and flag this issue.
% - The final matrix-monoid example is recorded as an observation, not
%   expanded into a proof: AB=CB with A != C.
% - Handwritten abbreviations and very faint show-through are not
%   reconstructed beyond what the mathematical lines support.

\subsection*{Seminar 4 - Page 2}

\noindent\underline{Exercise 6.}

\noindent\textbf{(a)} Show that every semigroup $S$ can be embedded in a
monoid: there exist a monoid $M$ and an injective semigroup
homomorphism
\[
  f:S\longrightarrow M, \qquad f(xy)=f(x)f(y).
\]

\noindent\textbf{(b)} There also exists a monoid that cannot be embedded
in a group: a monoid $M$ for which there are no group $G$ and
injective monoid homomorphism $f:M\to G$.

\medskip
\noindent\underline{Proof.}

\noindent\textbf{(a)} Let $S$ be a semigroup.

\smallskip
\noindent\textbf{I.} If $S$ is already a monoid, take $M=S$ and
\[
  i:S\longrightarrow M=S,\qquad i(x)=x.
\]

\noindent\textbf{II.} Suppose $S$ is a semigroup but not a monoid.
Take $\infty\notin S$ and put
\[
  M\overset{\mathrm{not}}{=}S\cup\{\infty\}.
\]
For $x,y\in M$ the operation is defined, as written in the notes, by
\[
\begin{gathered}
  x\circ y\ \overset{\mathrm{def}}{=}\ xy
     \quad(x,y\in S),\\
  \infty\circ x\ \overset{\mathrm{def}}{=}\
     x\circ\infty\ \overset{\mathrm{def}}{=}\ x
     \quad(x\in S).
\end{gathered}
\]
Then $(M,\circ,\infty)$ is a monoid. Consider
\[
  f:S\longrightarrow S\cup\{\infty\}=M,\qquad f(x)=x;
\]
thus $S$ is embedded in the monoid $M$.

\medskip
\noindent\textbf{(b)} The notes consider the monoid
\[
  \mathbb Z_n=\{\widehat{0},\widehat{1},\ldots,
      \widehat{n-1}\},\qquad (\mathbb Z_n,\cdot).
\]
Suppose that there exist a group $G$ and an injective monoid
homomorphism $f:\mathbb Z_n\to G$. Thus
\[
  f(xy)=f(x)f(y),\qquad f(\widehat{1})=e.
\]
The handwritten cancellation argument is
\[
\begin{aligned}
  f(\widehat{3})f(\widehat{2})
    &=f(\widehat{1})f(\widehat{2})\\
  &\Longrightarrow f(\widehat{3})=f(\widehat{1})\\
  &\overset{f\,\mathrm{injective}}{\Longrightarrow}
      \widehat{3}=\widehat{1}\qquad\text{(false).}
\end{aligned}
\]

\noindent Another example written at the bottom is the matrix monoid
$(M_n(K),\cdot)$, with matrices $A,B,C\in M_n(K)$ such that
\[
  AB=CB,\qquad A\neq C.
\]
; no preamble or forced page break.
% EDITORIAL NOTES:
% - Exercise 6 has two parts. The proof distinguishes an existing monoid
%   from a semigroup without an identity, then adjoins a new symbol, infinity.
% - The source explicitly writes x*infinity = infinity*x = x. It does not
%   separately write infinity*infinity = infinity; that case is required
%   to make the new operation a monoid operation.
% - The source first writes Z_n generically. The subsequent equality
%   3-hat*2-hat = 1-hat*2-hat, with 3-hat != 1-hat, works for n=4;
%   it is NOT true for arbitrary n. We retain Z_n and flag this issue.
% - The final matrix-monoid example is recorded as an observation, not
%   expanded into a proof: AB=CB with A != C.
% - Handwritten abbreviations and very faint show-through are not
%   reconstructed beyond what the mathematical lines support.

\subsection*{Seminar 4 - Page 2}

\noindent\underline{Exercise 6.}

\noindent\textbf{(a)} Show that every semigroup $S$ can be embedded in a
monoid: there exist a monoid $M$ and an injective semigroup
homomorphism
\[
  f:S\longrightarrow M, \qquad f(xy)=f(x)f(y).
\]

\noindent\textbf{(b)} There also exists a monoid that cannot be embedded
in a group: a monoid $M$ for which there are no group $G$ and
injective monoid homomorphism $f:M\to G$.

\medskip
\noindent\underline{Proof.}

\noindent\textbf{(a)} Let $S$ be a semigroup.

\smallskip
\noindent\textbf{I.} If $S$ is already a monoid, take $M=S$ and
\[
  i:S\longrightarrow M=S,\qquad i(x)=x.
\]

\noindent\textbf{II.} Suppose $S$ is a semigroup but not a monoid.
Take $\infty\notin S$ and put
\[
  M\overset{\mathrm{not}}{=}S\cup\{\infty\}.
\]
For $x,y\in M$ the operation is defined, as written in the notes, by
\[
\begin{gathered}
  x\circ y\ \overset{\mathrm{def}}{=}\ xy
     \quad(x,y\in S),\\
  \infty\circ x\ \overset{\mathrm{def}}{=}\
     x\circ\infty\ \overset{\mathrm{def}}{=}\ x
     \quad(x\in S).
\end{gathered}
\]
Then $(M,\circ,\infty)$ is a monoid. Consider
\[
  f:S\longrightarrow S\cup\{\infty\}=M,\qquad f(x)=x;
\]
thus $S$ is embedded in the monoid $M$.

\medskip
\noindent\textbf{(b)} The notes consider the monoid
\[
  \mathbb Z_n=\{\widehat{0},\widehat{1},\ldots,
      \widehat{n-1}\},\qquad (\mathbb Z_n,\cdot).
\]
Suppose that there exist a group $G$ and an injective monoid
homomorphism $f:\mathbb Z_n\to G$. Thus
\[
  f(xy)=f(x)f(y),\qquad f(\widehat{1})=e.
\]
The handwritten cancellation argument is
\[
\begin{aligned}
  f(\widehat{3})f(\widehat{2})
    &=f(\widehat{1})f(\widehat{2})\\
  &\Longrightarrow f(\widehat{3})=f(\widehat{1})\\
  &\overset{f\,\mathrm{injective}}{\Longrightarrow}
      \widehat{3}=\widehat{1}\qquad\text{(false).}
\end{aligned}
\]

\noindent Another example written at the bottom is the matrix monoid
$(M_n(K),\cdot)$, with matrices $A,B,C\in M_n(K)$ such that
\[
  AB=CB,\qquad A\neq C.
\]
; no preamble or forced page break.
% EDITORIAL NOTES:
% - Exercise 6 has two parts. The proof distinguishes an existing monoid
%   from a semigroup without an identity, then adjoins a new symbol, infinity.
% - The source explicitly writes x*infinity = infinity*x = x. It does not
%   separately write infinity*infinity = infinity; that case is required
%   to make the new operation a monoid operation.
% - The source first writes Z_n generically. The subsequent equality
%   3-hat*2-hat = 1-hat*2-hat, with 3-hat != 1-hat, works for n=4;
%   it is NOT true for arbitrary n. We retain Z_n and flag this issue.
% - The final matrix-monoid example is recorded as an observation, not
%   expanded into a proof: AB=CB with A != C.
% - Handwritten abbreviations and very faint show-through are not
%   reconstructed beyond what the mathematical lines support.

\subsection*{Seminar 4 - Page 2}

\noindent\underline{Exercise 6.}

\noindent\textbf{(a)} Show that every semigroup $S$ can be embedded in a
monoid: there exist a monoid $M$ and an injective semigroup
homomorphism
\[
  f:S\longrightarrow M, \qquad f(xy)=f(x)f(y).
\]

\noindent\textbf{(b)} There also exists a monoid that cannot be embedded
in a group: a monoid $M$ for which there are no group $G$ and
injective monoid homomorphism $f:M\to G$.

\medskip
\noindent\underline{Proof.}

\noindent\textbf{(a)} Let $S$ be a semigroup.

\smallskip
\noindent\textbf{I.} If $S$ is already a monoid, take $M=S$ and
\[
  i:S\longrightarrow M=S,\qquad i(x)=x.
\]

\noindent\textbf{II.} Suppose $S$ is a semigroup but not a monoid.
Take $\infty\notin S$ and put
\[
  M\overset{\mathrm{not}}{=}S\cup\{\infty\}.
\]
For $x,y\in M$ the operation is defined, as written in the notes, by
\[
\begin{gathered}
  x\circ y\ \overset{\mathrm{def}}{=}\ xy
     \quad(x,y\in S),\\
  \infty\circ x\ \overset{\mathrm{def}}{=}\
     x\circ\infty\ \overset{\mathrm{def}}{=}\ x
     \quad(x\in S).
\end{gathered}
\]
Then $(M,\circ,\infty)$ is a monoid. Consider
\[
  f:S\longrightarrow S\cup\{\infty\}=M,\qquad f(x)=x;
\]
thus $S$ is embedded in the monoid $M$.

\medskip
\noindent\textbf{(b)} The notes consider the monoid
\[
  \mathbb Z_n=\{\widehat{0},\widehat{1},\ldots,
      \widehat{n-1}\},\qquad (\mathbb Z_n,\cdot).
\]
Suppose that there exist a group $G$ and an injective monoid
homomorphism $f:\mathbb Z_n\to G$. Thus
\[
  f(xy)=f(x)f(y),\qquad f(\widehat{1})=e.
\]
The handwritten cancellation argument is
\[
\begin{aligned}
  f(\widehat{3})f(\widehat{2})
    &=f(\widehat{1})f(\widehat{2})\\
  &\Longrightarrow f(\widehat{3})=f(\widehat{1})\\
  &\overset{f\,\mathrm{injective}}{\Longrightarrow}
      \widehat{3}=\widehat{1}\qquad\text{(false).}
\end{aligned}
\]

\noindent Another example written at the bottom is the matrix monoid
$(M_n(K),\cdot)$, with matrices $A,B,C\in M_n(K)$ such that
\[
  AB=CB,\qquad A\neq C.
\]
